% -*- coding: utf-8 -*-
% !TEX program = xelatex
\documentclass[plain,noanswer,medmath]{../../template/jnuexam}
\usepackage{../../template/kysx}

\begin{document}


\examtitle{1999}{2}


\loadproblems{1}{1999P1}%载入试卷一


\makepart{填空题}{本题共5小题，每小题3分，满分15分}


\begin{problem}
曲线$\begin{cases}
  x = \e^t\sin 2t \\
  y = \e^t\cos t \\
\end{cases}$在点$(0,1)$处的法线方程为 \fillin{}．
\end{problem}


\begin{problem}
设函数$y = y(x)$由方程$\ln \left(x^2 + y \right) = x^3 y + \sin x$确定，
则$\left. \frac{\dy}{\dx} \right|_{x = 0} =$ \fillin{}．
\end{problem}


\begin{problem}
$\int \frac{x + 5}{x^2 - 6x + 13} \dx =$ \fillin{}．
\end{problem}


\begin{problem}
函数$y = \frac{x^2}{\sqrt{- x^2}}$在区间$\left[\frac{1}{2},\frac{\sqrt{3}}{2} \right]$
上的平均值为 \fillin{}．
\end{problem}


\useproblem{1}{1}{3}%【同试卷一第一(3)题】


\makepart{选择题}{本题共5小题，每小题3分，满分15分}


\useproblem{1}{2}{2}%【同试卷一第二(2)题】


\begin{problem}
设$\alpha (x) = \int_0^{5x} \frac{\sin t}{t}\dt$，
$\beta (x) = \int_0^{\sin x} \left(1 + t \right)^{1/t} \dt$，
则当$x \to 0$时$\alpha (x)$是$\beta (x)$的\pickout{}
\begin{abcd}
\item 高阶无穷小．
\item 低阶无穷小．
\item 同阶但不等价的无穷小．
\item 等价无穷小．
\end{abcd}
\end{problem}


\useproblem{1}{2}{1}%【同试卷一第二(1)题】


\begin{problem}
“对任意给定的$\varepsilon \in \left(0, 1 \right)$，总存在正整数$N$，当$n \ge N$时，
恒有$\left| x_n - a \right| \le 2\varepsilon$”是数列$\left\{x_n \right\}$收敛于$a$的\pickout{}
\begin{abcd}
\item 充分条件但非必要条件．
\item 必要条件但非充分条件．
\item 充分必要条件．
\item 既非充分条件又非必要条件．
\end{abcd}
\end{problem}


\begin{problem}
记行列式$\begin{vmatrix}
   x - 2 &  x - 1 &  x - 2 & x - 3 \\
  2x - 2 & 2x - 1 & 2x - 2 & 2x - 3 \\
  3x - 3 & 3x - 2 & 4x - 5 & 3x - 5 \\
      4x & 4x - 3 & 5x - 7 & 4x - 3
\end{vmatrix}$为$f(x)$，则方程$f(x) = 0$的根的个数为\pickout{}
\begin{abcd}
\item $1$．
\item $2$．
\item $3$．
\item $4$．
\end{abcd}
\end{problem}


\makepart{}{本题满分5分}

\begin{problem}
求$\lim_{x \to 0} \frac{\sqrt{1 + \tan x} - \sqrt{1 + \sin x}}{x\ln \left(1 + x \right) - x^2}$．
\end{problem}


\makepart{}{本题满分6分}

\begin{problem}
计算$\int_1^{+\infty} \frac{\arctan x}{x^2}\dx$．
\end{problem}


\makepart{}{本题满分7分}

\begin{problem}
求初值问题$\begin{cases}
  \left(y + \sqrt{x^2 + y^2} \right)\dx - x\dy = 0\enspace(x > 0) \\
  y\big|_{x = 1} = 0 \\
\end{cases}$的解．
\end{problem}


\makepart{}{本题满分7分}%【同试卷一第七题】（分值不同）

\useproblem{1}{7}{0}


\makepart{}{本题满分8分}

\begin{problem}
已知函数$y = \frac{x^3}{(x - 1)^2}$，求
\begin{enumerate}
\item 函数的增减区间及极值；
\item 函数图形的凹凸区间及拐点；
\item 函数图形的渐近线．
\end{enumerate}
\end{problem}


\makepart{}{本题满分8分}

\begin{problem}
设函数$f(x)$在闭区间$\left[- 1,1 \right]$上具有三阶连续导数，
且$f(- 1) = 0$，$f(1) = 1$，$f'(0) = 0$，证明：
在开区间$\left(- 1,1 \right)$内至少存在一点$\xi$，使$f'''\left(\xi \right) = 3$．
\end{problem}


\makepart{}{本题满分9分}%【同试卷一第五题】（分值不同）

\useproblem{1}{5}{0}


\makepart{}{本题满分6分}

\begin{problem}
设$f(x)$是区间$\left[0, +\infty \right)$上单调减少且非负的连续函数，
$$a_n = \sum_{i = 1}^n f(k) - \int_1^n f(x)\dx\quad \left(n = 1,2, \cdots \right),$$
证明数列$\left\{a_n \right\}$的极限存在．
\end{problem}


\makepart{}{本题满分8分}

\begin{problem}
设矩阵$A = \begin{pmatrix}
   1 &  1 & -1 \\
  -1 &  1 & 1 \\
   1 & -1 & 1
\end{pmatrix}$，矩阵$X$满足$A^* X = A^{-1} + 2X$，其中$A^*$是$A$的伴随矩阵，求矩阵$X$．
\end{problem}


\makepart{}{本题满分5分}

\begin{problem}
设向量组
$$\alpha_1=(1,1,1,3)^T, \alpha_2=(-1,-3,5,1)^T,
  \alpha_3=(3,2,-1,p+2)^T, \alpha_4=(-2,-6,10,p)^T.$$
\begin{enumerate}
\item $p$为何值时，该向量组线性无关？并在此时将向量$\alpha = \left(4,1,6,10 \right)^T$
      用$\alpha_1$, $\alpha_2$, $\alpha_3$, $\alpha_4$线性表示；
\item $p$为何值时，该向量组线性相关？并在此时求出它的秩和一个极大线性无关组．
\end{enumerate}
\end{problem}


\end{document}
